\section{Do the APPS rankings survive changes in source problems?}
\label{app:source-uncertainty}

The deployment ranges above answer: how much does a frozen panel's outcome
vary across populations with a fixed attack rate and budget? A second
question is whether a ranking depends on which APPS problems enter that
panel. We address it with a secondary analysis designed after the reported
rankings were known. We do not change the trained monitors, primary cell,
or original point estimates.

Each of 1,000 outer bootstrap draws samples 744 source problem IDs with
replacement, keeping each problem's honest and working-backdoor solution
together. Every monitor receives the same source draw. Within each draw,
the existing evaluator averages 256 deployment populations, each with 15
attacks, 1,473 honest submissions, and 29 audits. It uses shared population
draws and random tie ordering across monitors. AUROC is computed on the
balanced outer panel. The table gives paired percentile intervals over
source draws, not the central outcome ranges of individual deployment
populations. Nested Monte Carlo makes the loss intervals approximate.

\begin{table}[H]
\centering
\small
\setlength{\tabcolsep}{4pt}
\caption{Secondary source-problem bootstrap, conditional on the frozen fits
and split. Every difference is left minus right. Positive AUROC favors the
left monitor; negative violations favor it. Intervals are marginal 95\%
percentile intervals, not simultaneous confidence guarantees.}
\label{tab:apps-source-uncertainty}
\begin{tabular}{@{}L{0.48\linewidth}rr@{}}
\toprule
Model and comparison & $\Delta$AUROC interval & $\Delta$violations interval \\
\midrule
Gemma: detailed prompt $-$ SAE & $[0.005,0.043]$ & $[-1.18,3.92]$ \\
Gemma: verdict $-$ neutral residual & $[0.014,0.034]$ & $[-3.73,-1.78]$ \\
Qwen3.5: basic prompt $-$ text & $[-0.014,0.029]$ & $[1.50,4.61]$ \\
Qwen3.5: verdict $-$ neutral residual & $[0.030,0.053]$ & $[-5.48,-3.33]$ \\
\bottomrule
\end{tabular}
\end{table}

The result separates two conclusions. The verdict-adjacent residual's
advantage over the neutral residual persists under source resampling on both
models. The two prediction--action inversions are less settled. Gemma's
detailed prompt has a supported AUROC advantage over its SAE probe, but the
action-loss difference can change sign. Qwen3.5's basic prompt has a supported
action-loss disadvantage against text, but its small AUROC advantage can
change sign. Thus neither pair establishes both sides of a population-level
ranking reversal at this uncertainty level. They remain observed reversals
on the locked panel, not universal monitor orderings.

These intervals hold the fitted observers and model checkpoints fixed. They
do not include uncertainty from refitting, choosing layers or penalties,
selecting the contrasts after seeing outcomes, or changing the attacker.
Independent training splits and adaptive attacks require additional tests.
The analysis reuses the existing deployment evaluator and records the score
file hashes, seed, resampling unit, and replicate counts in its result files.
